% LSPV-3 REPORT::RETSOURCE1::2026-10-05T06:39:00-07:00
\documentclass[11pt]{article}
\usepackage[margin=.88in]{geometry}
\usepackage{amsmath,amssymb}
\begin{document}
\section*{LSPV-3 RETSOURCE1: source stop for SAT-semantic return}
\textbf{STATUS: OPEN.} P versus NP remains open. This report proves a source/interface statement only. It proves no SAT lower bound.

\section*{CORPUS REVIEWED}
Refetched and obeyed AGENTS.md, CODEX.md, the pnp-research skill, and root main.tex. Surveyed current tasks and repository search surfaces; checked the newest verified LSPV-2 predecessor DUALSHORT1-R2, PREDTYPE1, SOURCESTOP1, RETURNTRACE1/NORMBLOW1 surfaces, LSPV-4 RETURNTYPE1, RETURNMAP2 candidate attacks, current MAIN surfaces, and HEC/CRS-PPC/CCSPC source-stop lineage. Primary literature was rechecked in the 2025 van Melkebeek--Mocelin Sdroievski full paper and the Shaltiel--Umans AM reconstruction line.

\section*{Exact source target}
The checked vMS reconstruction is parameterized by an externally chosen function f used to instantiate the targeted HSG. In the applications, f maps its input to a truth table of a chosen hard language L at a shorter length. When derandomization fails, the reconstructor is run to compute a requested bit
\[
f(z)_w=L(w).
\]
In the all-input simulation, Merlin supplies a candidate distinguishing input and Arthur runs the reconstructor; after fixing Arthur randomness, the underlying computation has parallel SAT-oracle access. The acceptance test is whether the reconstructed bit equals one.

Thus the source-certified semantic output is a bit of the externally selected hard function f, not an ordinary CNF Q satisfying
\[
F\in SAT\iff Q\in SAT
\]
for an arbitrary original SAT instance F.

\section*{PROVED/EXACT relative to checked source}
Within the checked vMS/SU reconstruction interface, no source-defined NP-typed disjunction, failure certificate, or fixed-coin predicate was found whose truth is both equivalent to satisfiability of an arbitrary original SAT instance and free of unresolved current-scale use of the assumed SAT decider. The source reconstruction theorem guarantees recovery of f from a distinguisher under its hypotheses. It does not state a SAT self-reduction or SAT-semantic return map.

Therefore the Boundary-Squeeze return relation
\[
R_{\rm ret}(F,\mathcal S,Q),\qquad F\in SAT\iff Q\in SAT
\]
is not supplied by the checked reconstruction source. This is a source stop, not an impossibility theorem about every future construction.

\section*{Quantifier guard}
The missing theorem must uniformly map every relevant SAT input F and allowed assumed SAT decider A to a returned ordinary SAT instance Q, preserve satisfiability, preserve SAME_REALIZATION, have sublinear return exponent, and charge construction, acquisition, setup, classification, writing, workspace, and every call to A. The existing reconstruction theorem instead fixes a targeted generator based on f and reconstructs a requested coordinate of f from a distinguisher. These quantifiers and semantics are not interchangeable.

\section*{Conservative charge separation}
PREDTYPE1 remains controlling. An auxiliary NP acquisition predicate may be compiled to SAT, but its SAT instance preserves that auxiliary predicate, not original-SAT satisfiability. If acquisition uses n^u queries of length n^beta to a decider of exponent K, charge
\[
u+\beta K
\]
plus deterministic costs. Only a separately proved semantic return edge may contribute
\[
\eta+\theta K
\]
to recursive closure. RETURNTYPE1 further requires an ordinary explicit-CNF or model-matched random-access representation; a compact evaluator silently changes the target language unless a reduction is proved.

\section*{RETURNMAP2 hostile audit}
Direct restriction to a smaller CNF requires a satisfying-assignment selector; arbitrary fixed restrictions fail, and SAT self-reduction reacquires the unknown exponent. Generic verifier-to-CNF repairs type only for the verifier actually encoded and does not supply the missing semantic biconditional. Same-representation self-modeling remains an uninstantiated theorem schema. Many-query return remains conditional accounting until source-certified ordinary SAT queries exist. None repairs the source stop.

\section*{Literature transfer audit}
The Shaltiel--Umans line develops resilient reconstruction for AM hitting-set generators. This is directly relevant to local reconstruction but does not provide SAT-semantic self-reduction. The vMS 2025 application reconstructs hard-language truth-table bits and uses SAT-oracle computation inside the reconstruction protocol. This supports the source stop. QBF strategy extraction and proof-complexity interpolation are structural analogies only unless an exact SAT-semantic reduction is supplied.

\section*{DUPLICATION / CONTRADICTION AUDIT}
PREDTYPE1 already separated acquisition from return. RETURNTRACE1 already reported no ordinary SAT return object. RETSOURCE1 adds the requested primary-source search and identifies the actual source semantic target, namely f(z)_w=L(w). It therefore strengthens rather than contradicts those reports. RETURNTYPE1 independently blocks a compact-evaluator substitution.

\section*{WHAT CHANGED RELATIVE TO ROOT MAIN}
Root main asked to search the exact reconstruction proof and, if no suitable NP-typed SAT-equivalent predicate exists, prove a source-stop lemma. That search has now been carried out on the checked vMS reconstruction source. The correct next target is no longer another return-exponent optimization inside the existing source. It is a new source-specific semantic-return theorem or a different architecture.

\section*{ADOPTED / REJECTED / RETIRED PRIOR COMPONENTS}
Adopt PREDTYPE1, RETURNTRACE1, NORMBLOW1, RETURNTYPE1, SOURCESTOP1, and strict charge separation. Retire as a current-source route the inference from resilient reconstruction of f to a smaller ordinary SAT instance preserving arbitrary original-SAT satisfiability. Retain RETURNMAP2 Candidate C and Candidate D only as OPEN theorem schemas, not instantiated source edges.

\section*{Barrier ledger}
RETSOURCE1 relativizes and is not a nonrelativizing hinge. An oracle can make reconstruction easier without creating the missing SAT-semantic biconditional. Algebraic access likewise does not repair semantic typing. Arbitrary-algorithm acquisition, representation-invariant no-sharing, nonrelativization, nonalgebrization, and the final contradiction remain OPEN.

\section*{Single next target}
Either construct and prove a new source-specific SAT-semantic return theorem with exact ordinary-CNF/model-matched output and theta below one, or stop Boundary-Squeeze on the current vMS reconstruction source and redirect LSPV-3 effort to the main source-matched acquisition/reconstruction spine. No further optimization of the current RMV return exponent is justified before such a theorem exists.

\section*{FINAL STATUS}
\textbf{PROVED/EXACT relative to the checked source:} vMS reconstruction recovers coordinates of the chosen hard function and does not itself supply the required arbitrary-SAT semantic return relation. \textbf{OPEN:} a new SAT-semantic return theorem, arbitrary-algorithm acquisition, nonrelativization, nonalgebrization, and P versus NP.

\begin{thebibliography}{3}
\bibitem{vMS25} N. Mocelin Sdroievski and D. van Melkebeek, \emph{Instance-Wise Hardness and Refutation versus Derandomization for Arthur--Merlin Protocols}, Computational Complexity 34:15, 2025.
\bibitem{SU09} R. Shaltiel and C. Umans, \emph{Low-End Uniform Hardness versus Randomness Tradeoffs for AM}, SIAM J. Comput. 39(3):1006--1037, 2009.
\bibitem{SU05} R. Shaltiel and C. Umans, \emph{Simple Extractors for All Min-Entropies and a New Pseudorandom Generator}, J. ACM 52(2):172--216, 2005.
\end{thebibliography}
% END_OF_VERIFIED_REPORT::LSPV-3::RETSOURCE1
\end{document}
